\section{Martingales} \label{Ch4:Sec:Martingales}

\Cref{Ch4:Thm:Chernoff} got its strength from a very particular structure: a sum of independent contributions, each of them small. Plenty of quantities we care about have no such structure---the chromatic number of a random graph is not a sum of anything---and concentrate just as sharply anyway. What they do have is that they can be revealed a little at a time, with each revelation moving the expectation only slightly, and that is enough.

Throughout this section, fix random variables $X_0, X_1, X_2, \ldots$ on some probability space $(\Omega, \mathcal{P}(\Omega), \Pr)$.

The sequence to have in mind is the one where $X_k$ is our best guess at some quantity after $k$ things have been revealed, so that each term is the average of the term after it.

\begin{boxdefinition}[Martingale] \label{Ch4:Def:Martingale}
    We say $X_0, X_1, \ldots$ is a \textbf{martingale} if the following all hold:
    \begin{enumerate}
        % [CORRECTED] was: "for all $n \in \N$" in this item and in item 3 --- $\N$ starts at $1$ everywhere else in
        % these notes (the Zykov sequence $(G_n)_{n \in \N}$ of 1.3.3 starts at $G_1$), which leaves the step from
        % $X_0$ to $X_1$ governed by nothing: $X_0 = 0$ with $X_k = 10^6$ for $k \geq 1$ would qualify. The tree
        % below averages the root from its two children, which is item 3 at $n = 0$.
        \item for all integers $n \geq 0$, $\abs{X_{n+1} - X_n} \leq 1$
        % [FILLED] the lecture's second condition was illegible in the raw notes; the requirement below is the one the third condition needs, since $\E\of{X \mid A}$ divides by $\Pr(A)$
        \item the condition below is imposed only where $\Pr\of{X_0 = x_0, \, \ldots, \, X_n = x_n} \neq 0$, the conditional expectation being undefined otherwise
        % [CORRECTED] was: "$\E\of{X_{i + 1} \mid X_0 = x_0, \, X_1 = x_1, \, \cdots, \, X_{n-1} = x_{n-1}} = x_n$" --- $i$ is
        % free, and conditioning only as far as $x_{n-1}$ leaves the right-hand side $x_n$ undetermined by the event
        % on the left, so the condition reads "for all $x_n \in \R$, some fixed number equals $x_n$" and no
        % nondegenerate sequence can satisfy it. The tree below averages each node from the two beneath it, which
        % is this condition with the history conditioned on in full.
        \item for all integers $n \geq 0$ and all $x_0, \ldots, x_n \in \R$,
        \begin{align*}
            \E\of{X_{n + 1} \mid X_0 = x_0, \, X_1 = x_1, \, \cdots, \, X_n = x_n} = x_n
        \end{align*}
    \end{enumerate}
\end{boxdefinition}

\subsection{Exposing a Random Graph}

Let's study the expected value of chromatic numbers (cf. \Cref{Ch2:Def:Chromatic_Number}). Consider the family of random graphs $G_{3, \frac{1}{2}}$ of graphs on $3$ vertices with the probability of an edge between them being $\frac{1}{2}$ (cf. \Cref{Ch2:Def:Random_Graph}). This makes $\chi\of{G_{3, \frac{1}{2}}}$ a random variable. What would we expect its expectation to be?

We can heuristically see that the expected chromatic number is two based on the following figure:

\begin{figure}[H]
    \centering
    \begin{tikzpicture}[every node/.style = {font = \small, inner sep = 2pt}]
        \node (R)   at ( 0.00, 6.3) {\exposedtriad{unexposed}{unexposed}{unexposed}{2}};
        \node (A)   at (-3.00, 4.2) {\exposedtriad{present}{unexposed}{unexposed}{\frac{9}{4}}};
        \node (B)   at ( 3.00, 4.2) {\exposedtriad{absent}{unexposed}{unexposed}{\frac{7}{4}}};
        \node (AA)  at (-4.50, 2.1) {\exposedtriad{present}{present}{unexposed}{\frac{5}{2}}};
        \node (AB)  at (-1.50, 2.1) {\exposedtriad{present}{absent}{unexposed}{2}};
        \node (BA)  at ( 1.50, 2.1) {\exposedtriad{absent}{present}{unexposed}{2}};
        \node (BB)  at ( 4.50, 2.1) {\exposedtriad{absent}{absent}{unexposed}{\frac{3}{2}}};
        \node (AAA) at (-5.25, 0.0) {\exposedtriad{present}{present}{present}{3}};
        \node (AAB) at (-3.75, 0.0) {\exposedtriad{present}{present}{absent}{2}};
        \node (ABA) at (-2.25, 0.0) {\exposedtriad{present}{absent}{present}{2}};
        \node (ABB) at (-0.75, 0.0) {\exposedtriad{present}{absent}{absent}{2}};
        \node (BAA) at ( 0.75, 0.0) {\exposedtriad{absent}{present}{present}{2}};
        \node (BAB) at ( 2.25, 0.0) {\exposedtriad{absent}{present}{absent}{2}};
        \node (BBA) at ( 3.75, 0.0) {\exposedtriad{absent}{absent}{present}{2}};
        \node (BBB) at ( 5.25, 0.0) {\exposedtriad{absent}{absent}{absent}{1}};
        \node[anchor = east, gray] at (-6.1, 6.3) {$X_0$};
        \node[anchor = east, gray] at (-6.1, 4.2) {$X_1$};
        \node[anchor = east, gray] at (-6.1, 2.1) {$X_2$};
        \node[anchor = east, gray] at (-6.1, 0.0) {$X_3$};
        \draw (R) -- (A) node[midway, above left, inner sep = 1pt] {$e_1 \in G$};
        \draw (R) -- (B) node[midway, above right, inner sep = 1pt] {$e_1 \notin G$};
        \draw (A) -- (AA);
        \draw (A) -- (AB);
        \draw (B) -- (BA);
        \draw (B) -- (BB);
        \draw (AA) -- (AAA);
        \draw (AA) -- (AAB);
        \draw (AB) -- (ABA);
        \draw (AB) -- (ABB);
        \draw (BA) -- (BAA);
        \draw (BA) -- (BAB);
        \draw (BB) -- (BBA);
        \draw (BB) -- (BBB);
    \end{tikzpicture}
\end{figure}

Fix an order $e_1, e_2, e_3$ on the three possible edges and look at them one at a time. An edge not yet looked at is drawn dashed, since it is there with probability $\frac{1}{2}$; one already looked at is drawn solid if it is there and rubbed out if it is not. So the root has all three dashed, its two children split on $e_1$, theirs on $e_2$, and the eight leaves are the eight graphs on three vertices.

Under each picture is the expected chromatic number of a graph from the family it stands for---at a leaf, just that graph's: $1$ for the empty graph, $2$ for one or two edges, $3$ for the triangle. Each is the average of the two below it, so the $2$ at the root is $\E\of{\chi\of{G_{3, \frac{1}{2}}}}$.

Write $G$ for $G_{3, \frac{1}{2}}$ and $e_1, \ldots, e_m$ for the $m = \binom{3}{2}$ possible edges.

Define the following martingale:
% [FILLED] the eight right-hand sides were \sorry. The indices on the left were renumbered so that the display defines the sequence the sentence after it names: the author's $Y_0$ was never conditioned on, and the author's $X_1 = \E(X)$ left the $X_0$ of that sentence undefined.
\begin{align*}
    X &= \chi(G) \\
    Y_1 &= \ind{e_1 \in G} \\
    & \,\,\, \vdots \\
    Y_m &= \ind{e_m \in G} \\
    X_0 &= \E(X) \\
    X_1 &= \E\of{X \mid Y_1} \\
    X_2 &= \E\of{X \mid Y_1, Y_2} \\
    & \,\,\, \vdots
\end{align*}
Then, $X_0, X_1, X_2, \ldots$ is a martingale, and it is exactly what the tree draws: $X_k$ is the number written under the picture at the node reached by exposing $e_1, \ldots, e_k$. Its increments are never more than $\frac{1}{2}$, so the first condition is met with room to spare; $X_0$ is the $2$ at the root, and $X_m = X$ is the chromatic number itself, every $X_k$ past that being $X$ over again with nothing left to expose.

% [FILLED] the sentence broke off at the conditional expectation; everything from here to the end of the section was supplied, the statement of Azuma's inequality as the directive asked and the corollary as what the whole construction was for
We can adapt this to the general $G_{n, \frac{1}{2}}$ case. Define $Y_i$ and $X_i$ as above. Then, we can compute the conditional expectation of $\chi\of{G_{n, \frac{1}{2}}}$ on observing edges incident on vertices $1, \ldots, k$: write $X$ for $\chi\of{G_{n, \frac{1}{2}}}$, let $Y_k$ record the edges from vertex $k$ back to $1, \ldots, k-1$, and set $X_k = \E\of{X \mid Y_1, \ldots, Y_k}$. Since $Y_1$ records nothing, $X_0 = X_1 = \E(X)$; and $X_n = X$, with $X_k = X$ beyond that as before.

We are now exposing a vertex at a time rather than an edge at a time, which is what bounds the increments. Two graphs differing only at one vertex have chromatic numbers differing by at most $1$: give that vertex a color of its own. Pairing each unexposed configuration with itself under a different value of $Y_{k+1}$, the possible values of $X_{k+1}$ differ pairwise by at most $1$, and $X_k$ is their average, so $\abs{X_{k+1} - X_k} \leq 1$.

\subsection{Azuma's Inequality}

That condition is what buys us the following: a martingale cannot get far from where it started, and by the same margin \Cref{Ch4:Thm:Chernoff} gives a $\pm 1$ walk.

\begin{boxtheorem}[Azuma's Inequality] \label{Ch4:Thm:Azuma}
    Let $X_0, X_1, \ldots$ be a martingale and let $m$ be a positive integer. Then for every $\lambda > 0$,
    \begin{align*}
        \Pr\of{X_m - X_0 \geq \lambda \sqrt{m}} \leq e^{-\frac{\lambda^2}{2}}
    \end{align*}
\end{boxtheorem}

We won't prove it: it is the argument behind \Cref{Ch4:Thm:Chernoff}, with the chord bound applied to each increment conditionally on what has been exposed---which is what the bound on the increments is there to license.

So $\chi\of{G_{n, \frac{1}{2}}}$ sits in a window of width about $\sqrt{n}$ around its mean---much narrower than $\chi$ itself, which grows like $\frac{n}{\log n}$ (not something we will prove).

\begin{boxcorollary} \label{Ch4:Cor:Chromatic_Concentration}
    For every $n \geq 1$ and every $\lambda > 0$,
    \begin{align*}
        \Pr\of{\abs{\chi\of{G_{n, \frac{1}{2}}} - \E\of{\chi\of{G_{n, \frac{1}{2}}}}} \geq \lambda \sqrt{n}} \leq 2 e^{-\frac{\lambda^2}{2}}
    \end{align*}
\end{boxcorollary}
\begin{proof}
    The martingale above has $X_0 = \E\of{\chi\of{G_{n, \frac{1}{2}}}}$ and $X_n = \chi\of{G_{n, \frac{1}{2}}}$, so \Cref{Ch4:Thm:Azuma} at $m = n$ bounds $\Pr\of{X_n - X_0 \geq \lambda \sqrt{n}}$ by $e^{-\frac{\lambda^2}{2}}$. Applying it to $-X_0, -X_1, \ldots$, a martingale for the same reasons, bounds the other tail.
\end{proof}

\subsection{The Doob Martingale}

It turns out that the martingale we constructed for the chromatic number problem is a specific instance of a more general martingale, known as the Doob Martingale.

\begin{boxdefinition}[Doob Martingale] \label{Ch4:Def:Doob}
    For random variables $X, Y_1, Y_2, \ldots$ on some probability space, define
    % [CORRECTED] was: "$X_0 := X$" --- a Doob martingale runs from no information to all of it, so $X$ belongs at
    % the far end and the near end is the average over everything. With $X_0 = X$ the third condition of
    % \Cref{Ch4:Def:Martingale} fails already at $n = 0$: for $X$ uniform on $\set{0, 1}$ and $Y_1$ independent of
    % it, $\E\of{X_1 \mid X_0 = 0} = \frac{1}{2} \neq 0$. The construction of 4.3.1, which this one generalizes,
    % sets $X_0 = \E(X)$, and that is the $i = 0$ case of the display beside it with nothing conditioned on.
    % ("4.3.1" there is the subsection, Exposing a Random Graph, not the definition referenced two lines up.)
    \begin{align*}
        X_0 := \E\of{X}
        \qquad
        \forall i \geq 1, \, X_i := \E\of{X \mid Y_1, \ldots, Y_i}
    \end{align*}
    % [CORRECTED] was: "It is possible to show that $X_0, X_1, X_2, \ldots$ is a martingale, called the \textbf{Doob
    % Martingale}." --- \Cref{Ch4:Def:Martingale} makes $\abs{X_{n+1} - X_n} \leq 1$ part of the word "martingale",
    % and that is not free here: $X = 10 Y_1$ with $Y_1 = \pm 1$ gives $\abs{X_1 - X_0} = 10$. The averaging
    % condition does hold for every $X$ and every $Y_i$; the increment bound is a hypothesis, which is what the
    % sentence after this box and the $1$-coordinatewise-Lipschitz condition below both go on to supply.
    It is possible to show that $X_0, X_1, X_2, \ldots$ always satisfies the averaging condition of \Cref{Ch4:Def:Martingale}; once its increments are bounded by $1$ it is therefore a martingale, called the \textbf{Doob Martingale}.
\end{boxdefinition}

\subsection{McDiarmid's Inequality}

What makes \Cref{Ch4:Def:Doob} worth having is that its increments are controlled as soon as the $Y_i$ are independent and $X$ depends on no one of them too heavily.

\begin{boxdefinition}[Coordinatewise Lipschitz]
    Given metric spaces $\parenth{X_i, d_i}_{i=1}^{n}$ and $(Y, u)$, we say a function $f : X_1 \times \cdots \times X_n  \to Y$ is \textbf{$k$-coordinatewise-Lipschitz} if $u\of{f\of{\mathbf{a}}, f\of{\mathbf{b}}} \leq k$ if $\mathbf{a}$ and $\mathbf{b}$ differ in exactly one coordinate.
\end{boxdefinition}

Such an $f$, fed through the Doob martingale of its arguments, concentrates as sharply as a sum would.

\begin{boxproposition}[McDiarmid's Inequality] \label{Ch4:Prop:McDiarmid}
    Suppose $Y_1, \ldots, Y_n$ are independent random variables on domains $\mathcal{Y}_1, \ldots, \mathcal{Y}_n$. Suppose $f : \mathcal{Y}_1 \times \cdots \times \mathcal{Y}_n \to \R$ is $1$-coordinatewise-Lipschitz. Denoting $X = f\of{Y_1, \ldots, Y_n}$, we have that
    \begin{align*}
        \Pr\of{X \geq \E\of{X} + a} \leq e^{-\frac{a^2}{2n}}
    \end{align*}
    % [CORRECTED] was: "for all $a \in \R$" --- at $a < 0$ the bound is below $1$ while the probability tends to
    % $1$, and it fails for every admissible $f$: chaining the Lipschitz condition through $n$ single-coordinate
    % changes gives $X \geq \E(X) - n$ always, so at $a = -n$ the left-hand side is $1$ and the right-hand side is
    % $e^{-n/2}$. \Cref{Ch4:Thm:Chernoff} carries the same restriction, and Azuma's $\lambda > 0$ is where it
    % comes from.
    for all $a \geq 0$.
\end{boxproposition}
\begin{proof}
    Consider the Doob Martingale $X_i = \E\of{X \mid Y_1, \ldots, Y_i}$. Suppose we've fixed $y_1, \ldots, y_{i - 1}$. Define
    \begin{align*}
        g_i(y) = \E\of{f\of{y_1, \ldots, y_{i-1}, y, Y_{i+1}, \ldots, Y_n}}
    \end{align*}
    \sorry % [CLAUDE] there is no need to fill this \sorry for now: it is a placeholder. At some point, Wes will post notes about Azuma and McDiarmid, which I will give you to integrate into these notes (WITH ATTRIBUTION).
\end{proof}

One application of \Cref{Ch4:Prop:McDiarmid} is something called the Smiley Face Lemma, from a paper of Alan Frieze and Wes~\cite{FriezePegdenSubadditive}.

% [CORRECTED] was: "$D \ssq \R$" --- $D$ is the codomain of the embedding and the copy is said to lie in $\R^2$,
% so $D$ is a region of the plane. That is the whole of the reason. An earlier draft of this marker added that
% $\ssq \R$ makes the lemma false rather than merely vacuous, which implied the corrected version is not false;
% the [SUSPECT] below says why it is.
% [SUSPECT] the lemma below is false for every bounded $D$, and asking for nonempty interior does not save it.
% Take $D$ the open unit disc: the $n$ points are spread over a square of area $n$, so the number landing in $D$ is
% Binomial$\parenth{n, \frac{\pi}{n}}$ and stays $\BigO{\log n}$ with high probability, giving $o(n)$ copies however
% small $\delta$ is. The same count kills any $D$ of finite area. What the statement seems to want is a copy of $S$
% inside \textit{some} translate of $D$, anywhere in the square, rather than an embedding into one fixed planar
% set --- and, for the count to be defined at all, a copy made of the sampled points rather than of arbitrary
% points of $D$. Left as written: this is the cited paper's statement to settle, not ours.
For a finite set $S$, any $\eps > 0$, and any $D \ssq \R^2$, an $\eps$-$D$ copy of $S$ in $\R^2$ is some embedding of $S$ into $D$ such that distinct points can be separated by $\eps$-balls.

Scattering $n$ points at random over a square of area $n$ leaves linearly many such copies.

\begin{boxlemma}[Smiley Face Lemma, \cite{FriezePegdenSubadditive}]
    For any fixed $S, \eps, D$ as above, there is some $\delta > 0$ such that $n$ random points in a $\sqrt{n} \times \sqrt{n}$ square have $\delta n$ many $\eps$-$D$ copies of $S$ (with high probability).
\end{boxlemma}
